\section{偏好记忆}\label{sec:pref}

\subsection{偏好是如何学到的}

偏好记忆有三个信息来源，它们在后台自动工作，不需要你做任何额外操作。\textbf{工具
执行结果}是最主要的来源：助手统计你在同类任务中对各工具的使用成败，当某个工具在
某类任务上表现出稳定优势（样本充足、成功率高、占比过半）时，生成「这类任务优先用
这个工具」的偏好；反之若某工具连续失败，其权重会被降低但不会直接删除，给它留出
恢复的余地。\textbf{你的明确指令}置信度最高：你说「以后都用简洁格式」「构建前先跑
测试」这类话时，会被立即记录为偏好。\textbf{行为纠正}同样有效：如果你反复否决助手
的某种做法，系统会记录「避免」偏好并削弱与之冲突的旧习惯。

每条偏好都带有置信度，并按类别归档：工具选择、输出风格、安全策略、工作流、操作
习惯与通用偏好。安全敏感操作相关的倾向会单独归入安全策略类，供审批环节参考。

\subsection{版本、回滚与跨场景}

偏好的每次变化都产生一个新版本，历史版本完整保留。这意味着两件事：其一，你可以
随时要求助手「把 xx 偏好回滚到上周的值」，回滚以追加新版本的方式完成，历史不被
破坏；其二，你可以要求「看看这条偏好的变化过程」，得到完整的时间线。

偏好还可以绑定\textbf{场景}。例如你在终端场景偏好简短输出，在文档撰写场景偏好详尽
输出——两者作为同一偏好的不同场景版本并存。当助手在某个场景下没有完全匹配的偏好
时，会参考属性最相近场景下已验证的偏好作为起点。在指定时间点回溯「当时生效的是
哪个版本」也是支持的，便于排查「行为什么时候变的」这类问题。

\subsection{管理你的偏好}

你可以用自然语言管理偏好：问「你现在记住了我哪些偏好」列出全部；说「不要再自动用
xx 工具了」删除或降级对应偏好；说「这条偏好在 yy 场景下不适用」调整场景绑定。
偏好的读取入口（注入对话上下文）\todo{偏好注入提示词接线完成后，补充「验证偏好
生效」的操作说明与截图}。
